\scanpage{84}
% archive foliation: 83
d'où
\[
\Gamma=\operatorname{Gal}(\overline{\mathbb Q}/\mathbb Q)
\longrightarrow\widetilde{\mathcal T}_\nu
\]
L'hom correspondant
\[
\Gamma\longrightarrow\unsure{\mathfrak S|\mathcal S}_{\nu+1}
\]
se fait via
\[
\Gamma\longrightarrow\mu_{\nu+1}(\mathbb C^{*})\simeq\mathbb Z/\nu+1
\xrightarrow{\mathrm{can}}\unsure{\mathfrak S|\mathcal S}_{\nu+1}\ldots
\]
L'intérêt de prendre $I$ comme
dessus est d'avoir une opération
de \inserted{$\mathbb Z/2\cdot$}$\mu_{n\,\mathbb Q}$ \deleted{\unsure{trans.}} sur la situation
$(X_0/I_0/\mathbb Q)$ donc de \inserted{$\mathbb Z/2\cdot$}$\mu_n(\overline{\mathbb Q})\ (=D_n)$
sur $X/I/\mathbb Q$ --- en fait, le
produit $\tfrac12$ direct
$(\Gamma\times\deleted{\unsure{\mathbb Z/2}})\underbrace{\mu_n(\overline{\mathbb Q})}_{\subset\overline{\mathbb Q}^{*}}$\unsure{$=$}
opère sur $X/I/\overline{\mathbb Q}$ donc
\deleted{s'envoie dans}
\[
(\Gamma\times\mathbb Z/2)\cdot\mu_n(\overline{\mathbb Q})
\unsure{\longrightarrow|\Longrightarrow}
\widetilde{\mathcal T}_\nu
\]
où l'hom $\mu_n(\overline{\mathbb Q})\longrightarrow\widetilde{\mathcal T}_\nu$
\margin{$n=\nu+1$}
